\documentclass[12pt,a4paper]{article}

\usepackage{fontawesome5}
\providecommand{\faPencilAlt}{\faIcon{pencil-alt}}

\usepackage{styles/couleurs}
\usepackage{styles/mise_en_page}
\usepackage{styles/boites}
\usepackage{styles/maths}
\usepackage{styles/figures}

\usepackage{enumitem}
\usepackage{array}
\usepackage{tabularx}
\usepackage{amsmath}
\usepackage{amssymb}
\usepackage{graphicx}
\usepackage{tikz}
\usepackage{xparse}

\usetikzlibrary{calc,angles,quotes,arrows.meta}

\setlist[itemize]{leftmargin=1.4em,itemsep=0.15em,topsep=0.2em}
\setlist[enumerate]{leftmargin=*,itemsep=0.22em,topsep=0.25em}

\newcommand{\lignepointillee}{\noindent\dotfill\par\vspace{0.45em}}

\newcommand{\lignesreponse}[1]{%
\begin{tcolorbox}[
enhanced,colback=white,colframe=NeutreC!40,boxrule=0.5pt,
arc=1.5mm,left=3mm,right=3mm,top=2mm,bottom=1mm,
before skip=2pt,after skip=2pt]
\foreach \i in {1,...,#1}{\lignepointillee}
\end{tcolorbox}}

\newcommand{\titreseq}[1]{%
\begin{center}
{\sffamily\large\bfseries\color{SavoirC}#1}
\end{center}}

\begin{document}

\titreseq{Activité 4 -- Calculer le côté adjacent}

\section*{Partie A -- Le cosinus comme multiplicateur}

Dans le triangle \(ABC\) rectangle en \(B\), on donne :

\[
AC=12\text{ cm}
\qquad
\widehat{A}=38^\circ.
\]

\begin{center}
\begin{tikzpicture}[scale=0.9]
  \coordinate (A) at (0,0);
  \coordinate (B) at (6,0);
  \coordinate (C) at (6,4.7);
  \draw[line width=0.9pt] (A)--(B)--(C)--cycle;
  \draw (5.65,0)--(5.65,0.35)--(6,0.35);
  \draw pic[draw=ActionC,fill=ActionC!15,angle radius=9mm,"$38^\circ$"] {angle=B--A--C};
  \node[below left] at (A) {$A$};
  \node[below right] at (B) {$B$};
  \node[above right] at (C) {$C$};
  \node[above left] at ($(A)!0.52!(C)$) {$12$ cm};
  \node[below] at ($(A)!0.5!(B)$) {$?$};
\end{tikzpicture}
\end{center}

\begin{enumerate}
  \item Quelle est l'hypoténuse ?
  \item Quel est le côté adjacent à l'angle \(\widehat A\) ?
  \item Complète :
  \[
  \cos(\widehat A)=\ldots \div \ldots.
  \]
  \item À partir de cette égalité, complète :
  \[
  AB=
  \]
  \item Calcule \(AB\), arrondi au dixième de centimètre.
\end{enumerate}

\lignesreponse{7}

\remarque{\textbf{Aide — le sens d'un quotient}\\[0.3em]
Dire que $a\div b=q$, c'est dire que $q\times b=a$ :
le quotient est \textbf{le nombre qui, multiplié par $b$, donne $a$}.\\[0.3em]
\emph{Exemple : $15\div 5=3$ car $3\times 5=15$.}\\[0.3em]
Ici : $\cos(\widehat A)$ est donc le nombre qui, multiplié par $\ldots$,
donne $\ldots$}

\newpage

\section*{Partie B -- Une phrase opératoire}

Complète.

\begin{tcolorbox}[colback=SavoirC!5,colframe=SavoirC!60,boxrule=0.7pt,arc=2mm]
Pour calculer le côté adjacent connaissant l'hypoténuse et l'angle :

\[
\text{côté adjacent}
=
\]
\end{tcolorbox}

\section*{Partie C -- Calculer}

Pour chaque figure, commence par écrire une estimation et le nom du triangle rectangle.

\begin{enumerate}
  \item Dans le triangle \(DEF\) rectangle en \(E\),
        \(DF=8\) cm et \(\widehat D=24^\circ\). Calculer \(DE\).
  \item Dans le triangle \(GHI\) rectangle en \(H\),
        \(GI=15\) m et \(\widehat I=67^\circ\). Calculer \(HI\).
  \item Dans le triangle \(JKL\) rectangle en \(K\),
        \(JL=5{,}4\) cm et \(\widehat J=42^\circ\). Calculer \(JK\).
\end{enumerate}

\lignesreponse{15}

\section*{Partie D -- Contrôler}

Pour chacun des résultats précédents, vérifie :

\begin{itemize}
  \item que la longueur est positive ;
  \item qu'elle est plus petite que l'hypoténuse ;
  \item que l'unité et l'arrondi sont indiqués.
\end{itemize}

\end{document}
